\chapter[Electrónica y software de vuelo en detalle]{Electrónica, software de vuelo y estimación de estado en detalle}\label{cap:electronica}

Este capítulo baja a nivel de circuito y de código lo que el capítulo~\ref{cap:elec} planteó
en bloques. La electrónica propuesta cumple el anexo~\cite{anexo} con margen y la decisión
crítica, liberar el rotor al 80\,\% de la altitud máxima, resulta robusta frente a ruido,
valores atípicos y al pico de presión de la expulsión.

\begin{resultado}[Resultados principales]
\begin{itemize}[leftmargin=*]
  \item \textbf{Energía}: una 18650 de \SI{3000}{mAh} cubre \SI{2}{h} en la plataforma,
        el vuelo y \SI{1}{h} de recuperación con un factor de margen de 3,9.
  \item \textbf{Liberación}: con el filtro de Kalman y la lógica de votación, la liberación
        ocurre a \SI{-0.3}{m} (mediana) del 80\,\% real, sin disparos prematuros en
        \num{4000} vuelos simulados. Las lógicas ingenuas liberan antes de tiempo en el
        29--60\,\% de los vuelos si la expulsión produce un pico de presión.
  \item \textbf{Firmware}: se probó en la PC con los mismos datos de sensores (software en el
        lazo) y se recupera de un reinicio en cualquier estado. Un corte de \SI{0.3}{s}
        retrasa la liberación \SI{0.16}{s} como máximo, salvo que caiga en los \SI{0.3}{s}
        previos a la liberación: entonces el retraso llega a \SI{0.51}{s} ($\approx\SI{7}{m}$,
        79\,\% de $h_{\max}$).
  \item \textbf{Radio}: a \SI{915}{MHz} el enlace cierra a \SI{2}{km} con unos
        \SI{31}{dB} de margen (\SI{21}{dB} a \SI{2.4}{GHz}); tras el aterrizaje el enlace
        se pierde más allá de \SIrange{0.6}{1.1}{km}.
\end{itemize}
\end{resultado}

\begin{nota}[Alcance]
Todo lo que sigue es análisis y simulación; no hay ensayos. Los componentes marcados como
«verificado» tienen sus datos tomados de la hoja de datos citada; el resto son
«representativos» de su clase y deben medirse con el prototipo.
\end{nota}

%=====================================================================
\section{Arquitectura eléctrica}\label{sec:electronica-arq}

La arquitectura usa una sola celda 18650 y dos convertidores separados: un buck-boost de
\SI{3.3}{V} para la lógica y un elevador de \SI{5}{V} para el servo y las cámaras. Así,
los picos del servo no llegan al microcontrolador y la lógica funciona durante toda la
descarga de la celda (figura~\ref{fig:electronica-esq}).

\begin{figure}[H]
\centering
\begin{subfigure}{\textwidth}\centering
\begin{circuitikz}[font=\scriptsize, line width=0.5pt]
  \ctikzset{bipoles/length=0.8cm}
  % --- camino de potencia
  \draw (0,-2.6) to[battery1, l_={B1 18650}] (0,0)
        to[fuse, l={F1 \SI{3}{A}}] (1.3,0);
  \node[draw, fill=ffijo, minimum width=13mm, minimum height=8mm, align=center] (prot) at (2.15,0)
        {Protec.\\[-1pt]OV/UV/OC};
  \draw (1.3,0) -- (prot.west);
  \draw (prot.east) to[nos, l={S1 (\Req{E4})}] (4.1,0)
        to[R, l_={$R_s$ \SI{20}{m\ohm}}] (5.2,0) -- (9.4,0);
  \node[draw, fill=white, minimum width=13mm, minimum height=5mm] (ina) at (4.65,1.1) {INA219};
  \draw (4.1,0) |- ([yshift=-1mm]ina.west);
  \draw (5.2,0) |- ([yshift=-1mm]ina.east);
  \node[anchor=west, cgris] at (ina.east) {I$^2$C};
  \node[above, font=\scriptsize\bfseries] at (5.75,0) {VBAT};
  \draw (6.2,0) to[R, l_={\SI{1}{k\ohm}}, *-] (6.2,-1.3) to[leDo, l_={LED \Req{E6}}] (6.2,-2.6);
  \draw (8.0,0) to[R, l={\SI{100}{k\ohm}}, *-] (8.0,-1.3) to[R, l={\SI{100}{k\ohm}}] (8.0,-2.6);
  \draw (8.0,-1.3) to[short, *-] (7.3,-1.3) node[left] {ADC};
  \draw (0,-2.6) -- (10.6,-2.6);
  \draw (prot.south) -- (2.15,-2.6);
  % --- convertidores
  \node[draw, fill=ffijo, minimum width=20mm, minimum height=9mm, align=center] (u1) at (10.6,-1.0)
        {Buck-boost 3,3\,V\\[-1pt]TPS63020};
  \node[draw, fill=cfreno!10, minimum width=20mm, minimum height=9mm, align=center] (u2) at (10.6,1.0)
        {Elevador 5\,V\\[-1pt]2\,A (repr.)};
  \draw (9.4,0) to[short, *-] (9.4,1.0) -- (u2.west);
  \draw (9.4,0) -- (9.4,-1.0) -- (u1.west);
  \draw (u1.south) -- (10.6,-2.6);
  \draw (u2.south) -- ++(0,-0.2) node[ground]{};
  % --- riel de 3,3 V
  \draw (u1.east) -- (13.0,-1.0) node[right, align=left] {\textbf{+3V3}\\[-1pt]a la parte (b)};
  \draw (12.1,-1.0) to[C, l={$2\times\SI{22}{\micro F}$}, *-] (12.1,-1.8) node[ground]{};
  % --- riel de 5 V
  \draw (u2.east) -- (14.3,1.0) node[right, align=left] {a Q1 (servo)\\[-1pt]y Q2 (cám.)};
  \node[above, font=\scriptsize\bfseries] at (12.0,1.0) {+5V};
  \draw (12.4,1.0) to[eC, l_={\SI{470}{\micro F}}, *-] (12.4,-0.05) node[ground]{};
  \draw (13.7,1.0) to[D, l={TVS}, *-] (13.7,-0.05) node[ground]{};
  % --- baliza independiente
  \begin{scope}[shift={(11.6,-4.25)}]
    \draw[dashed, ccuerda] (0,0) rectangle (4.0,1.25);
    \draw (0.45,0.2) to[battery1, l_={B2}] (0.45,1.05) to[nos, l={S2 (\Req{E9})}] (2.1,1.05)
          to[buzzer] (3.6,1.05) -- (3.6,0.2) -- (0.45,0.2);
    \node[ccuerda, below, font=\tiny] at (2.0,0) {Baliza (\Req{C8}, \Req{E8}): sin conexión con el resto};
  \end{scope}
\end{circuitikz}
\caption{Potencia: celda, protección, interruptor, medición y rieles.}
\end{subfigure}

\vspace{4mm}
\begin{subfigure}{\textwidth}\centering
\begin{circuitikz}[font=\scriptsize, line width=0.5pt]
  \ctikzset{bipoles/length=0.7cm}
  \node[draw, fill=ffijo, minimum width=42mm, minimum height=62mm] (mcu) at (7.1,0) {};
  \node[align=center] at (6.7,-0.1) {\small\bfseries MCU\\ \footnotesize Cortex-M, 3,3\,V\\ \footnotesize lazo de 50\,Hz\\ \footnotesize perro guardián};
  \tikzset{per/.style={draw, fill=white, minimum width=30mm, minimum height=6mm, align=center}}
  % izquierda: bus I2C, SPI y RTC
  \node[per] (baro) at (1.6,2.6) {Barómetro MS5611};
  \node[per] (imu)  at (1.6,1.8) {IMU $\pm16\,g$ (repr.)};
  \node[per] (ina2) at (1.6,1.0) {INA219 (V, I)};
  \node[per] (fram) at (1.6,0.2) {FRAM MB85RC256V};
  \node[per] (sd)   at (1.6,-1.0) {Tarjeta SD};
  \draw[line width=1.2pt, cfijo] (3.7,2.6) -- (3.7,0.2);
  \foreach \n in {baro,imu,ina2,fram} \draw[cfijo] (\n.east) -- (3.7,0 |- \n.east) node[circ]{};
  \draw[cfijo, line width=1.2pt] (3.7,1.4) -- (mcu.west |- 0,1.4) node[right, black] {I2C};
  \draw[cfijo] (sd.east) -- (mcu.west |- sd.east) node[right, black] {SPI};
  \draw (mcu.west |- 0,-2.4) node[right] {VBAT\textsubscript{RTC}} -- (4.4,-2.4)
        to[battery1, l={CR1220 (\Req{F2})}] (3.0,-2.4) -- (2.6,-2.4) node[ground]{};
  % derecha
  \node[per] (gps)  at (14.1,2.6) {GPS (repr.)};
  \node[per] (xb)   at (14.1,1.8) {Radio XBee (\Req{X1})};
  \draw (gps.west) -- (mcu.east |- gps.west) node[left] {UART1};
  \draw (xb.west) -- (mcu.east |- xb.west) node[left] {UART2};
  \draw (xb.east) -- ++(0.3,0) node[antenna, anchor=south]{};
  \node[per, fill=cfreno!10] (srv) at (14.1,0.9) {Servo de traba};
  \draw (mcu.east |- srv.west) node[left] {PWM} -- (9.3,0.9) to[R, l={\SI{1}{k\ohm}}] (10.4,0.9) -- (srv.west);
  \draw (10.9,0.9) to[R, l_={\SI{10}{k\ohm}}, *-] (10.9,0.1) node[ground]{};
  \node[per, fill=frot] (hall) at (14.1,-0.5) {Hall + 2 imanes};
  \draw (hall.west) -- (12.3,-0.5) to[R, l={\SI{1}{k\ohm}}] (11.1,-0.5) -- (mcu.east |- 0,-0.5) node[left] {EXTI};
  \draw (12.3,-0.5) to[R, l={\SI{10}{k\ohm}}, *-] (12.3,0.2) node[right, font=\tiny] {+3V3};
  \draw (10.2,-0.5) to[C, l_={\SI{10}{nF}}, *-] (10.2,-1.25) node[ground]{};
  \node[per] (cmr) at (14.1,-2.0) {P-MOS Q1 (servo), Q2 (cám.)};
  \draw (cmr.west) -- (mcu.east |- cmr.west) node[left] {GPIO};
  \node[per] (sol) at (14.1,-2.8) {Paneles solares (\Req{E3})};
  \draw (sol.west) -- (mcu.east |- sol.west) node[left] {ADC (\Req{X7})};
\end{circuitikz}
\caption{Señales: buses, entradas con acondicionamiento y salidas.}
\end{subfigure}
\caption{Esquemático propuesto. «repr.»: componente representativo de su clase. La baliza
no comparte energía ni masa con el resto (\Req{C8}). Las resistencias de \SI{10}{k\ohm} en
el PWM y en el hall evitan movimientos del servo al arrancar y fijan el nivel del sensor de
salida a colector abierto.}
\label{fig:electronica-esq}
\end{figure}

\paragraph{Decisiones de diseño.}
\begin{itemize}
  \item \textbf{Celda}: 18650 de envase metálico (\Req{E1}, \Req{E2}) con lengüetas
        soldadas y cuna impresa con precinto. \Req{M4} prohíbe los contactos de resorte,
        así que el «portapila» de la tabla~\ref{tab:masa} no puede ser uno de resorte.
  \item \textbf{Buck-boost para la lógica}: con el servo bloqueado la tensión de la celda cae
        hasta \SI{2.8}{V} (sección~\ref{sec:electronica-pot}). Un regulador lineal de
        \SI{3.3}{V}, que necesita $\gtrsim\SI{3.5}{V}$ de entrada, podría reiniciar el
        procesador justo en la liberación si la celda está fría o poco cargada.
  \item \textbf{Barómetro a \SI{50}{Hz}}: el MS5611 necesita dos conversiones (presión y
        temperatura) de hasta \SI{9.04}{ms} cada una, es decir \SI{18.1}{ms}, por debajo de
        los \SI{20}{ms} del ciclo pero con poco margen. Si la lectura por I$^2$C no entra,
        la temperatura puede convertirse con OSR\,256 (\SI{0.6}{ms}), porque cambia
        lentamente.
  \item \textbf{Sensor hall}: va en la tapa superior fija, frente a dos imanes opuestos en
        el cubo (dos imanes mantienen el balance). No hace falta un anillo rozante. A
        \SI{1150}{rpm} da \SI{38}{Hz} (período de \SI{26}{ms}). El filtro RC tiene
        \SI{10}{\micro s} en el flanco de bajada y \SI{0.11}{ms} en el de subida (a través
        del pull-up de \SI{10}{k\ohm}); ambos son despreciables frente al período.
  \item \textbf{Hora de misión}: un RTC con pila de botón (permitida por \Req{E2}) sigue
        contando durante reinicios y cortes breves, como exige \Req{F2}.
\end{itemize}

\begin{table}[H]
\centering\small
\caption{Componentes y su función. V: datos verificados en la hoja de datos; R:
representativo.}\label{tab:electronica-comp}
\begin{tabularx}{\textwidth}{@{}llYc@{}}
\toprule
\textbf{Función} & \textbf{Componente} & \textbf{Dato clave} & \\\midrule
Medición de V e I (\Req{SN3}, \Req{SN8}) & INA219~\cite{electronica-ina219} & 0--\SI{26}{V}, \num{12} bits, $\pm\SI{320}{mV}$; con \SI{20}{m\ohm}: paso de \SI{0.5}{mA} (pide \SI{0.01}{A}) & V\\
Riel de \SI{3.3}{V} & TPS63020~\cite{electronica-tps63020} & entrada \SIrange{1.8}{5.5}{V}, hasta 96\,\%, \SI{25}{\micro A} en reposo & V\\
Altitud (\Req{SN1}) & MS5611~\cite{electronica-ms5611} & \SI{0.012}{mbar} RMS a OSR\,4096 ($\approx\SI{0.10}{m}$); conversión \SI{8.22}{ms} típ., \SI{9.04}{ms} máx. & V\\
Persistencia (\Req{F1}, \Req{F5}) & MB85RC256V~\cite{electronica-mb85rc256v} & \SI{32}{KB}, $10^{12}$ escrituras por byte, I$^2$C a \SI{1}{MHz} & V\\
Radio \SI{915}{MHz} (\Req{X1}) & XBee-PRO 900HP~\cite{electronica-xbee900} & \SI{24}{dBm}; \SI{-101}{dBm} a \SI{200}{kbit/s}; TX \SI{215}{mA}, RX \SI{29}{mA} & V\\
Radio \SI{2.4}{GHz} (\Req{X1}) & XBee 3 PRO~\cite{electronica-xbee3} & \SI{19}{dBm}; \SI{-103}{dBm}; TX \SI{135}{mA}, RX \SI{17}{mA} & V\\
Aceleración y giro (\Req{SN5}) & IMU de $\pm\SI{16}{g}$ & sesgo residual $\sim\SI{0.1}{m/s^2}$ tras calibrar & R\\
Posición (\Req{SN4}) & GPS de 1--\SI{10}{Hz} & $\approx\SI{30}{mA}$ en seguimiento & R\\
Traba & Servo MG90S & reposo \SI{10}{mA}, movimiento \SI{200}{mA}, bloqueo $\approx\SI{0.7}{A}$ (sin hoja de datos formal) & R\\
\bottomrule
\end{tabularx}
\end{table}

\begin{figure}[H]
\centering
\begin{tikzpicture}[x=1mm, y=0.34mm, font=\scriptsize]
  % cuerpo (escala vertical 0,34 mm/mm)
  \fill[ffijo] (0,0) rectangle (30,205);
  \draw[cfijo, line width=0.7pt] (0,0) rectangle (30,205);
  \foreach \ya/\yb/\yc/\txt in {0/40/20/{Baliza, lastre, cámara nadir},
      40/70/55/{Anillo de traba y servo}, 70/140/105/{Batería 18650},
      140/195/164/{Pila de PCB (MCU, sensores, radio)}, 195/205/204/{Tapa sup.: cámara cenit, hall}}{
    \draw[cfijo!60] (0,\ya) -- (30,\ya);
    \node[anchor=east, align=right, text width=40mm] at (-2,\yc) {\txt\\[-1pt]\textcolor{cgris}{\ya--\yb\,mm}};
  }
  \draw[rot] (8,205) rectangle (22,222);
  \node[anchor=west] at (34,213.5) {Cubo del rotor};
  \draw[cont] (8,222) rectangle (22,250);
  \node[anchor=west] at (34,236) {Copa del drogue};
  % palas plegadas (etapa 1)
  \fill[crot, opacity=0.35] (30,57) rectangle (32.5,213);
  \node[anchor=west, crot, font=\tiny, align=left] at (33,130) {palas\\[-1pt]plegadas\\[-1pt]57--213\,mm\\[-1pt](etapa 1)};
  % batería y PCB
  \draw[draw=cgris, fill=cgris!25] (11,72) rectangle (19,138);
  \foreach \y in {145,160,175,190} \draw[cfijo, line width=1pt] (2,\y) -- (28,\y);
  % antena en la zona libre inferior
  \draw[ccont, line width=1.2pt, decorate, decoration={coil, aspect=0.5, segment length=1.6mm, amplitude=1.2mm}] (25,8) -- (25,38);
  \node[ccont, anchor=west, font=\tiny] at (33,22) {antena (zona libre $<\SI{57}{mm}$)};
  \draw[ccont] (32.5,22) -- (26.5,22);
  % arnés
  \draw[ccuerda, line width=0.9pt] (25,40) -- (25,145);
  \draw[cfreno, line width=0.9pt] (5,55) -- (5,145);
  \draw[cgris, line width=0.9pt] (15,140) -- (15,145);
  \draw[crot, line width=0.9pt] (20,190) -- (20,203);
  % conectores
  \foreach \x/\y/\n in {25/145/J6, 5/145/J2, 15/141/J1, 20/195/J3, 10/198/J4, 9/30/J5}
    \node[marca, minimum size=3.6mm] at (\x,\y) {};
  \foreach \x/\y/\n in {25/145/6, 5/145/2, 15/141/1, 20/195/3, 10/198/4, 9/30/5}
    \node[font=\tiny\bfseries, text=white] at (\x,\y) {\n};
  % leyenda
  \node[anchor=north west, align=left, font=\scriptsize, text width=55mm] at (60,250) {%
    \textbf{Conectores con traba} (vibración \Req{S4}, choque \Req{S5}):\\[2pt]
    \textbf{1} batería: lengüetas soldadas + conector de 2 vías con traba (\Req{M4})\\
    \textbf{2} servo: 3 vías (5\,V, GND, PWM), cable trenzado\\
    \textbf{3} hall: 3 vías (3,3\,V, GND, señal), cable corto\\
    \textbf{4} cámara cenit: 5\,V, GND, control\\
    \textbf{5} cámara nadir: ídem; baliza con su propio cableado\\
    \textbf{6} antena: coaxial fino con conector a rosca o pegado\\[3pt]
    Interruptores S1 y S2 en orificios $\le\SI{10}{mm}$ (\Req{E5}) ubicados en los huecos
    entre palas plegadas. Todo cable con alivio de tensión y adhesivo en la ficha.};
\end{tikzpicture}
\caption{Arnés y conectores sobre la distribución de alturas de la
tabla~\ref{tab:largo} (escala vertical comprimida). En la etapa~1 las palas de carbono,
conductoras, cubren el cuerpo entre 57 y \SI{213}{mm}; la antena debe quedar debajo.}
\label{fig:electronica-arnes}
\end{figure}

%=====================================================================
\section{Presupuesto de potencia y caída de tensión}\label{sec:electronica-pot}

La energía no es crítica: la misión completa consume \SI{2.2}{Wh} de los \SI{8.6}{Wh}
utilizables. Lo crítico es la caída de tensión cuando arranca el servo, que define la
arquitectura de reguladores. La tabla~\ref{tab:electronica-pot} refina la
tabla~\ref{tab:energia} por estado; las potencias son estimaciones referidas a la entrada de
los convertidores (rendimientos supuestos de 0,90 y 0,88).

\begin{table}[H]
\centering\small
\caption{Potencia y energía por estado (\texttt{analysis/electronica\_potencia.py}). Riel de
\SI{3.3}{V}: \SI{173}{mA}, de los cuales la radio aporta \SI{32.7}{mA} en promedio (RX
continuo y \SI{20}{ms} de TX por paquete).}\label{tab:electronica-pot}
\begin{tabular}{@{}lcccr@{}}
\toprule
\textbf{Estado} & \textbf{Cámaras} & \textbf{Duración} & \textbf{Potencia} [W] & \textbf{Energía} [mWh]\\\midrule
\texttt{LAUNCH\_PAD} (\Req{E7}) & no & \SI{2}{h} & 0,69 & \num{1385}\\
\texttt{ASCENT} y \texttt{APOGEE} & sí & \SI{13}{s} & 2,74 & 10\\
\texttt{DESCENT} & sí & \SI{11}{s} & 2,74 & 8\\
\texttt{PROBE\_RELEASE} (servo \SI{0.3}{s}) & sí & \SI{3}{s} & 2,85 & 2\\
\texttt{PAYLOAD\_RELEASE} & sí & \SI{90}{s} & 2,74 & 68\\
\texttt{LANDED} (recuperación) & \SI{60}{s} & \SI{1}{h} & 0,73 & 726\\\midrule
\textbf{Total} & & & & \textbf{\num{2200}}\\
\bottomrule
\end{tabular}
\end{table}

Con \SI{3000}{mAh} a \SI{3.6}{V} y una reserva del 20\,\% por frío, envejecimiento y corte
a \SI{3.3}{V}, quedan \SI{8.64}{Wh}: un factor de margen de 3,9. La autonomía es de
\SI{12.5}{h} en el modo de plataforma y de \SI{3.2}{h} con todo encendido. Encender las
cámaras al detectar el despegue es lo que permite cumplir \Req{E7} con holgura.

\paragraph{Caída de tensión.} La carga de los convertidores es de potencia constante $P$,
así que la tensión en bornes $V$ de una celda con tensión en vacío $V_{oc}$ y resistencia
total $R$ (celda, protección, fusible, interruptor, shunt y cables) cumple $V = V_{oc} -
RP/V$:
\begin{equation}
  V = \frac{V_{oc} + \sqrt{V_{oc}^2 - 4RP}}{2}.
  \label{eq:electronica-caida}
\end{equation}
Con el servo bloqueado ($\approx\SI{0.7}{A}$ a \SI{5}{V}) y las cámaras encendidas,
$P=\SI{6.66}{W}$. Para $R=\SI{0.17}{\ohm}$ (celda nueva) y $V_{oc}=\SI{3.4}{V}$ (final de la
descarga) resulta $V=\SI{3.03}{V}$ con \SI{2.2}{A}; para $R=\SI{0.25}{\ohm}$ (celda fría o
envejecida), $V=\SI{2.81}{V}$. Ese es el peor caso: en una misión normal la celda llega a
la liberación con $\approx$87\,\% de carga ($V_{oc}\approx\SI{4.0}{V}$) y la tensión cae a
\SIrange{3.5}{3.7}{V}. Aun así, un regulador lineal quedaría al borde de su caída mínima,
mientras que el buck-boost sigue regulando hasta \SI{1.8}{V} (antes actúa el corte por baja
tensión de la protección).

\paragraph{Condensador del riel de 5 V.} Mientras el lazo del elevador reacciona
(tiempo $\Delta t$), el condensador entrega el escalón de corriente $I$ del servo:
\begin{equation}
  \Delta V \approx \frac{I\,\Delta t}{C} + I\,R_{\mathrm{ESR}}.
  \label{eq:electronica-cap}
\end{equation}
Con $I=\SI{0.7}{A}$, $R_{\mathrm{ESR}}=\SI{30}{m\ohm}$ y $C=\SI{470}{\micro F}$, la caída
es de \SIrange{95}{170}{mV} para $\Delta t=\SIrange{50}{100}{\micro s}$, un 2--3\,\% del
riel. Con \SI{220}{\micro F} llegaría a \SI{340}{mV}. El diodo TVS absorbe la fuerza
contraelectromotriz del motor al frenar.

%=====================================================================
\section{Estimación de altura y velocidad}\label{sec:electronica-kf}

Un filtro de Kalman~\cite{electronica-kalman1960} de dos estados fusiona la IMU y el
barómetro. Entrega la altura con \SI{0.06}{m} y la velocidad con \SI{0.07}{m/s} de desvío
teórico, sin el retardo de un promedio móvil.

\subsection{Modelo}
El estado es $\mathbf{x}=[h\;\;v]^{\mathsf T}$ (altura sobre el punto de \texttt{CAL} y
velocidad vertical, positivas hacia arriba). La IMU es la \emph{entrada}: con el eje $+Z$
hacia el nadir del anexo, la fuerza específica hacia arriba es $f=-a_Z$ y la aceleración
vertical es $u=f-g$. Con paso $\Delta t=\SI{0.02}{s}$:
\begin{equation}
  \mathbf{x}_k = \underbrace{\begin{bmatrix}1&\Delta t\\0&1\end{bmatrix}}_{F}\mathbf{x}_{k-1}
  + \underbrace{\begin{bmatrix}\Delta t^2/2\\ \Delta t\end{bmatrix}}_{G}(u_k + w_k),
  \qquad z_k = \underbrace{[1\;\;0]}_{H}\,\mathbf{x}_k + n_k,
  \label{eq:electronica-modelo}
\end{equation}
donde $z_k$ es la altura barométrica (atmósfera ISA referida a $p_0$),
$w_k\sim\mathcal N(0,\sigma_a^2)$ agrupa el error de la aceleración y
$n_k\sim\mathcal N(0,\sigma_b^2)$ el del barómetro. Las ecuaciones del filtro son las
habituales~\cite{electronica-simon2006}:
\begin{equation}
\begin{aligned}
  \hat{\mathbf x}^- &= F\hat{\mathbf x} + G u, & P^- &= FPF^{\mathsf T} + GG^{\mathsf T}\sigma_a^2,\\
  \nu &= z - H\hat{\mathbf x}^-, & S &= HP^-H^{\mathsf T} + \sigma_b^2,\\
  K &= P^-H^{\mathsf T}S^{-1}, & \hat{\mathbf x} &= \hat{\mathbf x}^- + K\nu,\quad P=(I-KH)P^-.
\end{aligned}
\label{eq:electronica-kf}
\end{equation}

\paragraph{Compuerta de innovación.} Una medición se rechaza si $\nu^2/S>25$ (5 desvíos).
Si se rechazan 50 seguidas (\SI{1}{s}), el filtro se «re-ancla»: suma $\nu^2$ a $P_{00}$ y
acepta la medición. La compuerta elimina los valores atípicos y el pico de presión de la
expulsión; el re-anclaje evita que un filtro divergido quede ciego.

\paragraph{Choque.} Si pasaron más de \SI{4}{s} desde el despegue (el motor ya se apagó:
su empuje da $f\approx9\,g$) y $|f|>6\,g$, la muestra de la IMU se descarta
(se mantiene la anterior y se aumenta $\sigma_a$). El choque de la expulsión no se integra
como velocidad falsa, y queda registrado como un voto de apogeo
(sección~\ref{sec:electronica-logica}).

\subsection{Régimen permanente y sintonía}
En régimen permanente el filtro es un sistema de segundo orden con frecuencia propia y
amortiguamiento que dependen solo del cociente de ruidos. Con las densidades espectrales
$q=\sigma_a^2\Delta t$ y $r=\sigma_b^2\Delta t$, la solución de la ecuación de Riccati del
filtro continuo~\cite{electronica-kalmanbucy1961} es
\begin{equation}
  P_{hh}=\sqrt2\,q^{1/4}r^{3/4},\quad P_{hv}=\sqrt{qr},\quad P_{vv}=\sqrt2\,q^{3/4}r^{1/4},
  \qquad
  \omega_n=\Bigl(\frac qr\Bigr)^{1/4}=\sqrt{\frac{\sigma_a}{\sigma_b}},\quad \zeta=\frac{1}{\sqrt2}.
  \label{eq:electronica-riccati}
\end{equation}
La tabla~\ref{tab:electronica-sintonia} da los valores elegidos. El ruido del MS5611 en
la hoja de datos es de \SI{0.10}{m}; en vuelo se supone el triple por el flujo alrededor
de los orificios, los cambios de temperatura y la turbulencia. Con
$\sigma_a=\SI{0.5}{m/s^2}$ y $\sigma_b=\SI{0.3}{m}$ resulta $\omega_n=\SI{1.29}{rad/s}$
($\SI{0.21}{Hz}$), $\sigma_h=\SI{0.057}{m}$ y $\sigma_v=\SI{0.074}{m/s}$. La ecuación de
Riccati discreta da los mismos valores, con ganancias $K=[\num{0.0359};\
\SI{0.0327}{s^{-1}}]$. Subir $\sigma_a$ hace al filtro más rápido y más ruidoso; bajarlo lo
vuelve sensible al sesgo del acelerómetro.

\begin{table}[H]
\centering\small
\caption{Sintonía del filtro.}\label{tab:electronica-sintonia}
\begin{tabularx}{\textwidth}{@{}lcY@{}}
\toprule
\textbf{Parámetro} & \textbf{Valor} & \textbf{Origen}\\\midrule
$\sigma_b$ (barómetro) & \SI{0.3}{m} & 3 veces el RMS del MS5611 a OSR\,4096 (\SI{0.10}{m}); supuesto\\
$\sigma_a$ (planeo y descenso) & \SI{0.5}{m/s^2} & sesgo residual $\sim\SI{0.1}{m/s^2}$ más inclinación de \ang{10}: $g(1-\cos\theta)=\SI{0.15}{m/s^2}$\\
$\sigma_a$ (motor) & \SI{5}{m/s^2} & vibración y saturación de la IMU durante $\approx\SI{2.6}{s}$\\
$\sigma_a$ (modo simulación) & \SI{3}{m/s^2} & sin IMU: el filtro solo usa la presión de \texttt{SIMP}\\
\bottomrule
\end{tabularx}
\end{table}

\paragraph{Errores de escala y de cero.} El criterio del 80\,\% es un cociente, y eso lo
protege. Si la altura medida es $h_m=s\,h+\delta$ (error de escala $s$ por la temperatura
real del aire y deriva $\delta$ de la presión desde el \texttt{CAL}), la condición
$h_m=0{,}8\,h_{m,\max}$ se cumple en
\begin{equation}
  h = 0{,}8\,h_{\max} - 0{,}2\,\delta/s.
  \label{eq:electronica-cociente}
\end{equation}
El error de escala se cancela y solo el 20\,\% de la deriva pasa a la liberación: una deriva
de \SI{0.5}{hPa} en \SI{2}{h} de espera ($\approx\SI{4}{m}$) la corre menos de \SI{1}{m}.

%=====================================================================
\section{Simulación del vuelo y comparación de lógicas}\label{sec:electronica-sim}

El script \texttt{analysis/electronica\_kalman.py} simula el vuelo, los sensores y tres
lógicas de decisión. La lógica propuesta (C) libera al 80\,\% con un error mediano de
\SI{-0.3}{m}; las lógicas ingenuas fallan por los valores atípicos o por el pico de la
expulsión.

\paragraph{Perfil sintético.} Plataforma \SI{5}{s}; motor de \SI{1.6}{s} a
\SI{80}{m/s^2} netos; planeo balístico con arrastre $k=\SI{5.26e-4}{m^{-1}}$, que lleva al
apogeo nominal de \SI{700}{m} a los \SI{12.1}{s}; expulsión cerca del apogeo; \SI{0.5}{s}
de caída del cuerpo solo; drogue con $\Vd=\SI{13.4}{m/s}$; tras la liberación, transición a
$\Vr=\SI{6.4}{m/s}$ con constante de tiempo de \SI{1.5}{s}. Sensores: barómetro con ruido
blanco $\sigma_b$ y valores atípicos de \SIrange{5}{30}{m} en el 0,2\,\% de las muestras;
IMU con sesgo, ruido de \SI{0.05}{m/s^2}, vibración del motor de \SI{3}{m/s^2},
oscilación de hasta \ang{15} bajo el drogue, choque de expulsión y saturación a $16\,g$.

\paragraph{Escenarios.} En el escenario~1 solo hay ruido y valores atípicos. El
escenario~2 agrega dos efectos supuestos (no medidos): un pulso de sobrepresión de la carga
de expulsión de \SIrange{1}{4}{kPa} durante \SIrange{0.05}{0.3}{s} (lectura \SIrange{83}{333}{m}
más baja) y un error de presión dinámica $k_q v^2/2g$ con $k_q\le0{,}02$ durante el ascenso.
En el Monte Carlo (\num{2000} vuelos por escenario) varían el empuje, el arrastre (apogeos de
\SIrange{520}{900}{m} entre los percentiles 5 y 95), $\Vd\sim\mathcal N(13{,}4;\,0{,}7)$,
$\sigma_b\in[0{,}1;\,0{,}5]$\,m, el sesgo, el momento de la expulsión ($-1$ a $+2$\,s respecto
del apogeo) y el choque (3 a $30\,g$).

\begin{figure}[H]
\centering
\begin{tikzpicture}
\begin{axis}[name=ah, width=\textwidth, height=5.0cm, xmin=0, xmax=125, ymin=-20, ymax=760,
    ylabel={Altura [m]}, xticklabels={}, legend style={font=\scriptsize, at={(0.99,0.97)}, anchor=north east},
    legend cell align=left]
  \addplot[cgris!60, line width=0.4pt] table[x=t, y=zb] {analysis/data/electronica_kalman_vuelo.dat};
  \addplot[black, line width=1.2pt] table[x=t, y=h] {analysis/data/electronica_kalman_vuelo.dat};
  \addplot[cfijo, line width=0.7pt, dashed] table[x=t, y=hk] {analysis/data/electronica_kalman_vuelo.dat};
  \legend{barómetro, real, Kalman}
  \draw[cfreno, dashed] (17.58,-20) -- (17.58,760) node[pos=0.06, right, font=\scriptsize] {apogeo det.};
  \draw[cfreno, dashed] (28.40,-20) -- (28.40,760) node[pos=0.06, right, font=\scriptsize] {liberación};
  \draw[ccuerda, dashed] (113.7,-20) -- (113.7,760) node[pos=0.94, left, font=\scriptsize] {suelo};
  \draw[cgris, densely dotted] (0,561) -- (125,561) node[pos=0.55, above, font=\scriptsize] {$0{,}8\,h_{\max}=\SI{561}{m}$};
\end{axis}
\begin{axis}[at={(ah.south west)}, anchor=north west, yshift=-2mm, width=\textwidth, height=3.6cm,
    xmin=0, xmax=125, ymin=-0.6, ymax=0.6, xlabel={Tiempo [s]}, ylabel={Error},
    legend style={font=\scriptsize, at={(0.99,0.97)}, anchor=north east}, legend columns=2]
  \addplot[cfijo, line width=0.5pt] table[x=t, y=eh] {analysis/data/electronica_kalman_vuelo.dat};
  \addplot[crot, line width=0.5pt] table[x=t, y=ev] {analysis/data/electronica_kalman_vuelo.dat};
  \legend{$\hat h-h$ [m], $\hat v-v$ [m/s]}
\end{axis}
\end{tikzpicture}
\caption{Vuelo nominal en lazo cerrado (escenario 2). La lógica C detecta el apogeo
\SI{0.52}{s} después del real y libera a \SI{560.5}{m} (79,9\,\% de \SI{701.4}{m}), a
\SI{0.04}{s} del cruce real del 80\,\%. En el ascenso rápido el error llega a \SI{15}{m} y
\SI{9}{m/s} (fuera de escala) por el error de presión dinámica supuesto ($\le\SI{8}{m}$). La
compuerta lo agranda: rechaza esas lecturas y el re-anclaje corrige la altura pero no la
velocidad (sin compuerta el máximo es de \SI{9}{m}). Ese error desaparece antes del apogeo: cerca de él
es de \SI{0.11}{m} y \SI{0.31}{m/s}, y en todo el descenso queda por debajo de
\SI{0.15}{m} y \SI{0.14}{m/s}. El pico final es un artefacto: el modelo detiene el CanSat al
tocar el suelo sin que la IMU registre el impacto.}
\label{fig:electronica-vuelo}
\end{figure}

\begin{figure}[H]
\centering
\begin{tikzpicture}
\begin{axis}[width=0.5\textwidth, height=5.4cm, xmin=-2.5, xmax=4.5, ymin=560, ymax=712,
    xlabel={$t-t_{\mathrm{apogeo}}$ [s]}, ylabel={Altura [m]},
    legend style={font=\scriptsize, at={(0.02,0.03)}, anchor=south west}, legend cell align=left,
    legend columns=1]
  \addplot[cgris!60, line width=0.4pt] table[x=tr, y=zb] {analysis/data/electronica_kalman_apogeo.dat};
  \addplot[black, line width=1.1pt] table[x=tr, y=h] {analysis/data/electronica_kalman_apogeo.dat};
  \addplot[ccuerda, line width=0.8pt] table[x=tr, y=he] {analysis/data/electronica_kalman_apogeo.dat};
  \addplot[cfijo, line width=0.8pt, dashed] table[x=tr, y=hk] {analysis/data/electronica_kalman_apogeo.dat};
  \legend{barómetro, real, EMA (B), Kalman (C)}
  \node[font=\scriptsize, align=center, cfreno] at (axis cs:2.9,620) {pico de expulsión:\\lectura en \SI{491}{m}};
  \draw[cfreno, -{Latex[length=1.6mm]}] (axis cs:1.9,612) -- (axis cs:0.75,575);
\end{axis}
\end{tikzpicture}\hfill
\begin{tikzpicture}
\begin{axis}[width=0.5\textwidth, height=5.4cm, xmin=-2.5, xmax=4.5, ymin=-30, ymax=30,
    xlabel={$t-t_{\mathrm{apogeo}}$ [s]}, ylabel={Velocidad [m/s]},
    legend style={font=\scriptsize, at={(0.98,0.97)}, anchor=north east}, legend cell align=left]
  \addplot[black, line width=1.1pt] table[x=tr, y=v] {analysis/data/electronica_kalman_apogeo.dat};
  \addplot[ccuerda, line width=0.8pt] table[x=tr, y=ve] {analysis/data/electronica_kalman_apogeo.dat};
  \addplot[cfijo, line width=0.8pt, dashed] table[x=tr, y=vk] {analysis/data/electronica_kalman_apogeo.dat};
  \legend{real, EMA (B), Kalman (C)}
\end{axis}
\end{tikzpicture}
\caption{Detalle del apogeo y de la expulsión (a $+\SI{0.52}{s}$) del vuelo nominal. El pico
de \SI{2.5}{kPa} lleva la lectura cruda a \SI{491}{m}, por debajo del umbral de
\SI{561}{m}: la lógica A libera en ese instante a \SI{700}{m}. La EMA baja hasta
\SI{599}{m} y su velocidad hasta \SI{-299}{m/s} (fuera de escala). El filtro C rechaza las muestras del pico
y no se aparta de la altura real.}
\label{fig:electronica-apogeo}
\end{figure}

\begin{table}[H]
\centering\small
\caption{Monte Carlo de \num{2000} vuelos por escenario. Retardo: tiempo entre el apogeo real
y su detección. Error de liberación: altura real al liberar menos $0{,}8\,h_{\max}$.
Prematura: por encima del 85\,\% de $h_{\max}$; tardía: por debajo del 75\,\% o sin
liberar.}\label{tab:electronica-mc}
\setlength{\tabcolsep}{4.5pt}
\begin{tabular}{@{}clccccccc@{}}
\toprule
& & \multicolumn{2}{c}{\textbf{Retardo apogeo} [s]} & \textbf{Apogeo} & \multicolumn{4}{c}{\textbf{Liberación}}\\
\cmidrule(lr){3-4}\cmidrule(lr){6-9}
\textbf{Esc.} & \textbf{Lógica} & med. & p95 & falso [\%] & err. med. [m] & p95 $|$err.$|$ [m] & prem. [\%] & tardía [\%]\\\midrule
1 & A: crudo     & 0,34 & 0,52 & 37,5 & $+0{,}2$ & 393 & 0,2 & 28,4\\
1 & B: EMA       & 0,54 & 0,58 & 0,8  & $-3{,}4$ & 3,8 & 0,0 & 0,8\\
1 & C: Kalman    & 0,70 & 0,78 & 0,0  & $-0{,}3$ & 0,5 & 0,0 & 0,0\\\midrule
2 & A: crudo     & $-0{,}30$ & 0,50 & 45,6 & $+107$ & 389 & 60,3 & 24,1\\
2 & B: EMA       & 0,50 & 0,56 & 13,4 & $-3{,}4$ & 146 & 28,8 & 0,9\\
2 & C: Kalman    & 0,70 & 0,80 & 0,1  & $-0{,}3$ & 0,6 & 0,0 & 0,0\\
\bottomrule
\end{tabular}
\end{table}

\paragraph{Lectura de la tabla.}
\begin{itemize}
  \item \textbf{A} toma el máximo de las lecturas crudas: un valor atípico positivo en el
        ascenso infla $h_{\max}$, las muestras siguientes quedan por debajo y declara un
        apogeo falso (37,5\,\%). Con $h_{\max}$ congelado bajo, libera tarde (28,4\,\%).
  \item \textbf{B} filtra los atípicos, pero su retardo de $\approx\SI{0.24}{s}$ hace que
        libere \SI{3.4}{m} tarde, y el pulso de expulsión la hace liberar casi en el apogeo
        en el 28,8\,\% de los vuelos.
  \item \textbf{C} detecta el apogeo algo más tarde (exige dos votos), pero estima
        $h_{\max}$ con \SI{0.05}{m} de error mediano. Su error de liberación tiene un p95 de
        \SI{0.5}{m} (escenario 1) y \SI{0.6}{m} (escenario 2). Los 2 apogeos «falsos» del
        escenario 2 se adelantan \SI{0.5}{s} con el CanSat expulsado aún subiendo a
        \SI{5}{m/s}; liberan igual al 80,3 y al 80,7\,\%.
  \item En la etapa 1 los errores RMS de C son \SI{0.09}{m} y \SI{0.14}{m/s}, mayores que
        los teóricos (\SI{0.06}{m}, \SI{0.07}{m/s}) porque el sesgo y la oscilación no son
        ruido blanco.
\end{itemize}

\begin{nota}[Límites del modelo]
El barómetro simulado es ruido blanco más pulsos. No incluye el retardo de los orificios de
ventilación, el transitorio térmico al salir del cohete ni el cambio del flujo alrededor del
cuerpo cuando arranca el rotor. Este último puede dar un escalón de pocos metros justo
después de liberar; no afecta la decisión, que ya se tomó, pero sí la velocidad medida en
esos segundos. El pulso de expulsión es una hipótesis de estrés: debe medirse con un ensayo
de expulsión en tierra con el CanSat dentro del cohete.
\end{nota}

%=====================================================================
\section{Lógica robusta y recuperación tras reinicios}\label{sec:electronica-logica}

Cada decisión irreversible combina varias evidencias, una ventana de tiempo plausible y un
respaldo independiente del barómetro. El estado se guarda de forma que un reinicio no la
repita ni la pierda.

\paragraph{Apogeo (2 de 3 votos).} V1: $\hat v<\SI{-1}{m/s}$ durante \SI{0.2}{s}; V2:
$\hat h_{\max}-\hat h>\SI{3}{m}$; V3: choque de expulsión ($|f|>6\,g$). El apogeo se
acepta si hay dos votos, han pasado \SI{4}{s} desde el despegue (bloqueo) y el filtro
convergió ($\sigma_v<\SI{2}{m/s}$). Si a los \SI{40}{s} del despegue no se detectó, se
declara igual. La guarda de
convergencia apareció en la prueba del firmware: sin ella, un reinicio a mitad del ascenso
dejaba $\hat v=0$, la IMU llevaba $\hat v$ a valores negativos y se declaraba un apogeo
falso.

\paragraph{Liberación.} Se ordena cuando se cumplen tres condiciones durante
$T_p=\SI{0.2}{s}$: el estado es \texttt{DESCENT}, $\hat h+\hat v\,T_p\le0{,}8\,h_{\max}$ (el
término $\hat vT_p$ compensa la espera) y
$t-t_{\mathrm{apo}}\ge0{,}2\,h_{\max}/V_{\lim}$ con $V_{\lim}=\SI{40}{m/s}$ (ningún descenso
real recorre el 20\,\% más rápido). El \textbf{respaldo por tiempo} sale de la caída desde
el reposo con arrastre cuadrático, $\Delta h=(V^2/g)\ln\cosh(gt/V)\approx Vt-(V^2/g)\ln2$:
\begin{equation}
  t_{\mathrm{res}} = t_{\mathrm{apo}} + \frac{0{,}2\,h_{\max}}{V_{\min}} + \frac{V_{\min}}{g}\ln2,
  \qquad V_{\min}=\SI{10}{m/s}.
  \label{eq:electronica-timer}
\end{equation}
Para \SI{700}{m} da \SI{14.7}{s}, más tarde que la liberación barométrica nominal
(\SI{10.8}{s} después de detectar el apogeo); en el Monte Carlo nunca actuó. Si el
barómetro falla, el temporizador libera a $\approx$76, 74 o 71\,\% de $h_{\max}$ para
$\Vd=12$, 13,4 o \SI{15}{m/s}. Es un modo degradado: no cumple el 80\,\% de \Req{C4} al
pie de la letra (para \SI{700}{m} libera entre 27 y \SI{64}{m} más abajo), pero garantiza
que la segunda etapa se abra.

\paragraph{Salud del barómetro y GPS.} El barómetro se declara en falla tras \SI{0.5}{s} sin
lecturas válidas, o si en el descenso difiere en más de \SI{50}{m} de la altura GPS durante
\SI{3}{s}. En ese caso $h_{\max}$ se toma del GPS y solo actúa el temporizador.

\paragraph{Persistencia (\Req{F1}, \Req{F5}).} Un registro de \SI{40}{bytes} (estado,
traba, modo, $p_0$, $h_{\max}$, última altura, hora del apogeo y contador de paquetes) se
escribe en la FRAM en dos copias alternadas con número de secuencia y CRC-32. Si un corte
interrumpe una escritura, la otra copia sigue siendo válida. Se escribe en cada cambio de
estado, en cada comando y en cada paquete, y a \SI{5}{Hz} durante el ascenso. Con unas 10
escrituras por segundo, las $10^{12}$ escrituras de la FRAM no son una limitación.

\paragraph{Perro guardián.} El temporizador de \SI{1}{s} solo se reinicia si el lazo de
control corrió en los últimos \SI{100}{ms}. Un bloqueo del bus I$^2$C o de la SD termina en
un reinicio limpio, y la tabla~\ref{tab:electronica-recup} define qué pasa después. Al
arrancar, el PWM del servo se fija según la traba guardada \emph{antes} de encender su riel:
un reinicio en vuelo no vuelve a trabar ni mueve el anillo.

\paragraph{Peor momento para un reinicio.} Si el corte cae en los \SI{0.3}{s} previos a la
liberación, el firmware pierde los \SI{0.3}{s} sin energía y vuelve a exigir los
\SI{0.2}{s} de persistencia: la liberación se atrasa hasta \SI{0.51}{s}, unos \SI{7}{m} a
\SI{13.4}{m/s} (79\,\% de $h_{\max}$ en lugar de 80\,\%). En cualquier otro momento el atraso
es de \SI{0.16}{s} como máximo.

\begin{table}[H]
\centering\small
\caption{Recuperación tras un reinicio, según el estado guardado. Última columna: prueba del
firmware en la PC con un corte de \SI{0.3}{s} (sin reinicio: liberación a \SI{28.38}{s}, giro
confirmado a \SI{30.30}{s} y \texttt{LANDED} a \SI{124.8}{s}).}
\label{tab:electronica-recup}
\begin{tabularx}{\textwidth}{@{}lYc@{}}
\toprule
\textbf{Estado guardado} & \textbf{Acción al arrancar} & \textbf{Prueba (corte en $t$)}\\\midrule
\texttt{LAUNCH\_PAD} & Conserva $p_0$, contador y traba; espera el despegue & ---\\
\texttt{ASCENT} & Filtro con $h$ del barómetro y $v$ libre ($\sigma_v=\SI{100}{m/s}$); bloqueo de \SI{2}{s}; $h_{\max}$ guardado & 5,5 a \SI{17.3}{s}: libera a 28,36--\SI{28.38}{s}\\
\texttt{APOGEE}, \texttt{DESCENT} & Pasa a \texttt{DESCENT}; recupera $h_{\max}$ y la hora del apogeo (RTC) para el temporizador & 17,7 a \SI{28.0}{s}: libera a 28,38--\SI{28.54}{s}; 28,1 a \SI{28.35}{s}: libera a 28,64--\SI{28.89}{s}\\
\texttt{PROBE\_RELEASE} & Reafirma «liberado» y vuelve a esperar el giro & 29 y \SI{30}{s}: giro confirmado a 30,30 y \SI{30.86}{s}\\
\texttt{PAYLOAD\_RELEASE} & Mantiene; solo espera el aterrizaje & 60 y \SI{120}{s}: \texttt{LANDED} \SI{1.5}{s} más tarde\\
\texttt{LANDED} & Mantiene; apaga cámaras a los \SI{60}{s} & ---\\
\bottomrule
\end{tabularx}
\end{table}

%=====================================================================
\section{Software de vuelo}\label{sec:electronica-fw}

El archivo \texttt{analysis/electronica\_firmware.cpp} (\num{408} líneas) contiene la
lógica completa en C++ para Arduino/PlatformIO. El acceso al hardware queda en funciones
\texttt{hal\_*} que se implementan con las bibliotecas de cada componente. El lazo principal
corre el control a \SI{50}{Hz}, la telemetría a \SI{1}{Hz} y lee comandos en cada vuelta.

\begin{table}[H]
\centering\small
\caption{Comandos (\texttt{CMD,<TEAM\_ID>,\ldots}). Los opcionales están permitidos por el
anexo; \texttt{SIM} y \texttt{SIMP} siguen la guía de la competencia~\cite{cansat2025}.}
\label{tab:electronica-cmd}
\begin{tabularx}{\textwidth}{@{}lYl@{}}
\toprule
\textbf{Comando} & \textbf{Acción y protecciones} & \textbf{Eco}\\\midrule
\texttt{CX,ON|OFF} & Activa la telemetría; también se activa sola en \texttt{APOGEE} (\Req{X4}) & \texttt{CXON}\\
\texttt{ST,hh:mm:ss|GPS} & Ajusta el RTC (\Req{F3}) & \texttt{ST13:35:59}\\
\texttt{CAL} & $p_0$ = media exponencial de \SI{1}{s}, $h_{\max}=0$, contador a cero; \textbf{solo en tierra} & \texttt{CAL}\\
\texttt{MEC,LATCH,ON|OFF} & \texttt{ON} libera (siempre); \texttt{OFF} traba (solo en \texttt{LAUNCH\_PAD}) & \texttt{MECLATCHON}\\
\texttt{MEC,CAM,ON|OFF} & Cámaras (prueba en tierra) & \texttt{MECCAMON}\\
\texttt{PWR,ON|OFF} (opc.) & Agrega la tensión del panel solar (\Req{X7}) & \texttt{PWRON}\\
\texttt{SIM,ENABLE|ACTIVATE|DISABLE} (opc.) & Modo simulación, solo en plataforma; \texttt{MODE}=\texttt{S} & \texttt{SIMACTIVATE}\\
\texttt{SIMP,<Pa>} (opc.) & Presión simulada; reemplaza al barómetro & \texttt{SIMP101325}\\
\bottomrule
\end{tabularx}
\end{table}

Un comando con otro número de equipo o con argumentos inválidos se descarta sin eco. El
paquete agrega, después de las dos comas del anexo, \texttt{ROTOR\_RPM}, \texttt{LATCH} y la
tensión del panel. Este paquete lo generó el firmware en la prueba en la PC
(sección~\ref{sec:electronica-fw}), \SI{0.6}{s} después de liberar; los sensores son
simulados y la tensión, la corriente y el panel son valores fijos del banco:
\begin{lstlisting}[numbers=none, basicstyle=\ttfamily\scriptsize, breaklines=true]
1000,13:36:26,26,F,PROBE_RELEASE,552.6,21.5,94.9,3.9,0.74,0.0,0.0,4.0,0.12,-0.05,-12.07,13:36:27,672.6,-31.4201,-64.1888,9,PWRON,,158,1,2.35
\end{lstlisting}
El anexo pide \texttt{ACCEL} en «grados por segundo al cuadrado»; se transmite en
\si{m/s^2} y conviene confirmar la unidad con la organización. También pide la longitud en
«grados Oeste»: el firmware la envía con signo (negativa al oeste), y la convención también
debe confirmarse.

\paragraph{rpm por interrupción.} La interrupción guarda el período entre pulsos y descarta
los menores de \SI{4}{ms} (rebotes). El lazo toma la mediana de los tres últimos períodos,
que elimina un pulso perdido o uno espurio, y la suaviza. El rotor se declara detenido si no
llega un pulso en $\max(\SI{0.5}{s},\,3T)$, con tope de \SI{1.5}{s}. Con
\texttt{micros()} la resolución es de \SI{1}{\micro s} sobre \SI{26}{ms}.

\lstinputlisting[language=C++, firstline=133, lastline=156, firstnumber=133,
  caption={Medición de rpm (extracto de \texttt{electronica\_firmware.cpp}).},
  label={lst:electronica-rpm}]{analysis/electronica_firmware.cpp}

\lstinputlisting[language=C++, firstline=237, lastline=267, firstnumber=237,
  caption={Apogeo por votación y liberación (extracto).},
  label={lst:electronica-estados}]{analysis/electronica_firmware.cpp}

\paragraph{Prueba en la PC (software en el lazo).} El banco
\texttt{analysis/electronica\_sil.cpp} compila el firmware sin cambios. Reproduce a
\SI{1}{kHz} los sensores del vuelo nominal de la figura~\ref{fig:electronica-vuelo}, envía
comandos, genera los pulsos del hall con un rotor que arranca con $\tau=\SI{1.5}{s}$ e
inyecta reinicios. Los resultados:
\begin{itemize}
  \item Las transiciones coinciden con el modelo en Python: \texttt{ASCENT} a
        \SI{5.08}{s}, \texttt{APOGEE} a \SI{17.58}{s} y liberación a \SI{28.38}{s}
        (\SI{28.40}{s} en Python).
  \item \texttt{PAYLOAD\_RELEASE} llega \SI{1.92}{s} después de liberar. Como el arranque
        real puede tardar \SIrange{3}{6}{s} (capítulo~\ref{cap:dim}), la espera antes del
        reintento es de \SI{8}{s}, no de \SI{5}{s}.
  \item El rotor detenido en el suelo se informa \SI{0.5}{s} después; un \texttt{CAL} en
        vuelo y un comando con otro número de equipo se rechazan.
  \item \texttt{PACKET\_COUNT} sigue la cuenta a través de los reinicios. El perro guardián
        se atendió en cada vuelta del lazo (cada \SI{1}{ms} del banco); esto solo prueba que
        la lógica no se bloquea, porque el banco no modela la demora real de la SD, del
        I$^2$C ni de la radio.
  \item \texttt{LANDED} se declara a \SI{124.8}{s}, \SI{11}{s} después del contacto. Es el
        mismo artefacto de la figura~\ref{fig:electronica-vuelo}: la IMU sintética no
        registra el impacto y $\hat v$ tarda en volver a cero. Con datos reales se espera
        menos, pero debe verificarse.
\end{itemize}

%=====================================================================
\section{Enlace de radio}\label{sec:electronica-rf}

Las dos bandas de \Req{X1} cierran el enlace a \SI{2}{km} con margen amplio, pero
\SI{915}{MHz} da unos \SI{10}{dB} más. El punto débil no es la distancia: es la antena del
CanSat, la sombra de las palas de carbono y la pérdida del enlace tras el aterrizaje.

El margen se calcula con la fórmula de Friis~\cite{electronica-friis1946}:
\begin{equation}
  M = P_t + G_{\mathrm{CanSat}}(\theta) + G_{\mathrm{tierra}} - L - 20\log_{10}\frac{4\pi d f}{c} - S,
  \label{eq:electronica-friis}
\end{equation}
con $S$ la sensibilidad del receptor y $\theta$ el ángulo entre el eje del CanSat y la
dirección a la estación. La antena del CanSat se modela como un dipolo vertical,
$G(\theta)=1{,}64\,[\cos(\tfrac\pi2\cos\theta)/\sin\theta]^2$, con un piso de
\SI{-20}{dBi} en el nulo.

\begin{table}[H]
\centering\small
\caption{Presupuesto de enlace a \SI{2}{km} de distancia oblicua, con el CanSat a
\SI{560}{m} ($\theta=\ang{74}$). Antenas de tierra y pérdidas: valores
representativos.}\label{tab:electronica-enlace}
\begin{tabular}{@{}lrr@{}}
\toprule
\textbf{Término} & \textbf{\SI{915}{MHz}} (900HP) & \textbf{\SI{2.44}{GHz}} (XBee 3 PRO)\\\midrule
Potencia transmitida $P_t$ & \SI{24.0}{dBm} & \SI{19.0}{dBm}\\
Antena del CanSat a $\theta=\ang{74}$ & \SI{1.6}{dBi} & \SI{1.6}{dBi}\\
Antena directiva de tierra, apuntada & \SI{8.0}{dBi} & \SI{10.0}{dBi}\\
Cables (a bordo y en tierra) & \SI{-1.5}{dB} & \SI{-2.0}{dB}\\
Polarización (oscilación de $\pm\ang{20}$) & \SI{-1.0}{dB} & \SI{-1.0}{dB}\\
Cuerpo, cables y palas (supuesto) & \SI{-3.0}{dB} & \SI{-3.0}{dB}\\
Pérdida en espacio libre & \SI{-97.7}{dB} & \SI{-106.2}{dB}\\
Sensibilidad $S$ & \SI{-101}{dBm} (\SI{200}{kbit/s}) & \SI{-103}{dBm}\\\midrule
\textbf{Margen} & \textbf{\SI{31.4}{dB}} & \textbf{\SI{21.4}{dB}}\\
\bottomrule
\end{tabular}
\end{table}

\begin{figure}[H]
\centering
\begin{tikzpicture}
\begin{axis}[width=0.9\textwidth, height=6.0cm, xmin=0, xmax=2.5, ymin=-15, ymax=80,
    xlabel={Distancia horizontal a la estación [km]}, ylabel={Margen [dB]},
    legend style={font=\scriptsize, at={(0.99,0.97)}, anchor=north east}, legend columns=2,
    legend cell align=left]
  \addplot[cfijo, line width=1.0pt] table[x=x, y=m915a] {analysis/data/electronica_enlace.dat};
  \addplot[crot, line width=1.0pt] table[x=x, y=m2400a] {analysis/data/electronica_enlace.dat};
  \addplot[cfijo, line width=0.7pt, dashed] table[x=x, y=m915b] {analysis/data/electronica_enlace.dat};
  \addplot[crot, line width=0.7pt, dashed] table[x=x, y=m2400b] {analysis/data/electronica_enlace.dat};
  \addplot[cfreno, densely dotted, line width=0.8pt, domain=0:2.5] {10};
  \legend{\SI{915}{MHz}, \SI{2.4}{GHz}, \SI{915}{MHz} en el suelo, \SI{2.4}{GHz} en el suelo, margen de desvanecimiento}
\end{axis}
\end{tikzpicture}
\caption{Margen del enlace con el CanSat a \SI{560}{m} (Friis y diagrama del dipolo) y ya
aterrizado, con la antena a \SI{0.3}{m} del suelo (dos rayos con reflexión $-1$,
antena de tierra a \SI{1.5}{m}). La caída cerca de 0 es el nulo del dipolo sobre la
estación.}
\label{fig:electronica-enlace}
\end{figure}

\paragraph{Conclusiones del enlace.}
\begin{itemize}
  \item \textbf{En vuelo} el margen mínimo es de \SI{20.8}{dB} (\SI{915}{MHz}) y
        \SI{10.8}{dB} (\SI{2.4}{GHz}), cuando el CanSat pasa sobre la estación y esta queda en
        el nulo del dipolo. El alcance en espacio libre con \SI{10}{dB} de reserva es de
        \SI{25}{km} y \SI{7.9}{km}.
  \item \textbf{En el suelo}, la reflexión hace que la pérdida crezca con $d^4$ a partir del
        punto de quiebre~\cite{electronica-rappaport2002}. El margen de \SI{10}{dB} se pierde a
        \SIrange{0.60}{0.66}{km} y el enlace se corta a \SIrange{1.07}{1.17}{km}, casi igual
        en las dos bandas. El cálculo supone la antena vertical; si el CanSat queda acostado,
        la desadaptación de polarización resta todavía más. Por eso el último paquete con GPS
        antes del contacto y la baliza son los que permiten recuperar el CanSat.
  \item \textbf{Zona de Fresnel}: su radio en el punto medio de \SI{2}{km} es de \SI{12.8}{m}
        (\SI{915}{MHz}) y \SI{7.8}{m} (\SI{2.4}{GHz}). Con el CanSat bajo, la estación debe
        estar en un lugar alto o sobre un mástil.
  \item \textbf{Antena del CanSat}: debe ir en la zona inferior libre ($<\SI{57}{mm}$,
        figura~\ref{fig:electronica-arnes}). En la etapa 1 las palas de carbono plegadas
        rodean el cuerpo entre 57 y \SI{213}{mm} y la sombrearían. A \SI{2.4}{GHz} un
        monopolo de $\lambda/4$ (\SI{3.1}{cm}) entra en esa zona. A \SI{915}{MHz} no entra
        (\SI{8.2}{cm}): hace falta una hélice compacta, con algunos dB menos de eficiencia,
        que el margen de \SI{10}{dB} extra absorbe. No puede sobresalir, porque \Req{S3}
        limita el largo total.
  \item \textbf{Antena de tierra}: Yagi con polarización vertical, apuntada a mano. Una
        circular perdería \SI{3}{dB} frente a la lineal del CanSat sin ganar nada mientras
        este baje vertical.
  \item \textbf{Configuración}: el 900HP trabaja en \SIrange{902}{928}{MHz}; hay que
        limitar sus canales a \SIrange{915}{928}{MHz} (\Req{X1}) y verificar que el salto de
        frecuencia lo admita. NetID o PAN ID igual al número de equipo (\Req{X2}) y envío
        unicast a la dirección de la estación (\Req{X3}). La telemetría ocupa
        $\approx\SI{2}{kbit/s}$, así que sobra velocidad.
\end{itemize}

%=====================================================================
\section{Síntesis y trabajo pendiente}\label{sec:electronica-sintesis}

\begin{resultado}[Recomendaciones]
\begin{enumerate}[leftmargin=*]
  \item Buck-boost para la lógica y elevador separado con \SI{470}{\micro F} y TVS para el
        servo; servo energizado solo después de fijar su PWM.
  \item Filtro de Kalman con la IMU como entrada, compuerta de innovación, votación 2 de 3,
        anticipación $\hat vT_p$ y respaldo por tiempo (ecuación~\ref{eq:electronica-timer}).
  \item FRAM con doble copia y CRC, RTC con pila y perro guardián condicionado al lazo de
        control. \texttt{CAL} solo en tierra.
  \item Radio de \SI{915}{MHz} con canales limitados a \SIrange{915}{928}{MHz}, antena en la
        zona inferior y Yagi vertical en tierra.
\end{enumerate}
\end{resultado}

\paragraph{Ensayos pendientes.} (1) Caída de tensión real con el servo bloqueado y la celda
fría. (2) Ruido del barómetro dentro del cuerpo, con flujo, y pulso de presión de una carga
de expulsión real. (3) Reinicios forzados en banco durante una reproducción del perfil.
(4) Prueba de alcance a \SI{2}{km}, con la antena montada y las palas plegadas. (5) Ajuste
de la espera de giro con las rpm reales del ensayo E2 (capítulo~\ref{cap:riesgos}).
